\section{Introduction}\label{sec:intro}

\subsection{Geometry, not the force law, controls contact fidelity}\label{sec:intro-motivation}

A computational contact algorithm must answer two local questions about a query point $\xbf$ near a
body $X$. First, is $\xbf$ inside $X$, and by how much? This is the gap $\gn$. Second, which way is
out? This is the normal $\nbf$. We use the sign convention $\gn<0\Leftrightarrow$ penetration throughout.
Once these two quantities are produced, a mature library of enforcement schemes (penalty, augmented
Lagrangian, mortar, Nitsche) converts them into nodal forces \citep{wriggers2006computational,
laursen2002computational,kikuchi1988contact}. On smooth, convex bodies this division of labour is
benign. The gap and the normal are smooth functions of position, every scheme converges to the same
answer, and the discussion centres on constraint enforcement.

% code: solvers/contact/gap.py::evaluate_gap
On geometrically complex bodies the same division of labour breaks down. The gap and the normal are
no longer smooth, or even single-valued, so the \emph{geometric representation} from which they are read
becomes the dominant control on accuracy. A rock joint, a cusped cam lobe, a faceted aggregate, or a
self-similar fractal boundary all carry contact information at length scales an ambient field of fixed
capacity cannot store. The thesis of this paper is one sentence: \emph{on rough and non-convex
geometry the representation of the boundary, not the constraint-enforcement scheme, controls contact
fidelity.} We make this thesis precise and give it a method: contact read from a chart
\emph{transition map} on a learned coordinate atlas. We then measure it against the incumbent
representation on a verification ladder that runs from closed-form benchmarks to a real rock joint.

\subsection{The level set and its two intrinsic weaknesses}\label{sec:intro-levelset}

% code: solvers/contact/gap.py::evaluate_gap
The incumbent representation is the level set. For a body $A$ one stores a signed-distance function
$\phi_A:\R^d\to\R$, negative inside the body and increasingly represented by a neural network, and reads
the gap and the normal directly from the field. The gap is the field value at the query point, and the
outward normal is its normalized gradient:
\begin{equation}
  \gn(\xbf)=\phi_A(\xbf),\qquad
  \nbf(\xbf)=\frac{\grad\phi_A(\xbf)}{\norm{\grad\phi_A(\xbf)}},
  \label{eq:sdf-gap-intro}
\end{equation}
so a single automatic-differentiation call delivers both quantities at once
\citep{liu2020implicit,kolev2003level}. The economy of \eqref{eq:sdf-gap-intro} is the appeal of the
representation, and the source of its weaknesses. To make $\phi_A$ a true signed \emph{distance}, one
trains it with the eikonal penalty, which pushes the gradient toward unit norm everywhere:
\begin{equation}
  \mathcal{L}_{\mathrm{eik}}
  =\mathbb{E}_{\xbf\sim\Omega}\!\left[\big(\norm{\grad\phi_A(\xbf)}-1\big)^2\right],
  \label{eq:eikonal}
\end{equation}
where $\Omega$ is the sampling domain and the expectation averages the unit-gradient residual over it.
This penalty is combined with surface and normal-alignment losses on a boundary point cloud. We use the
level set throughout the paper \emph{only as the comparison baseline}. It is a fair and capable baseline,
not a strawman. On the geometry it can represent, it is exactly as accurate as the chart
(Section~\ref{sec:verif}).

The level set carries two weaknesses that are intrinsic to the representation rather than artefacts of
training, and both bear directly on contact:
\begin{enumerate}
  \item \textbf{Medial-axis non-smoothness.} The exact Euclidean signed distance is only $C^0$ on the
        medial axis, the locus equidistant from two or more boundary points, where the closest-point map
        is multivalued and $\grad\phi$ jumps. In a concavity the medial set reaches close to the
        surface, so the gradient normal of \eqref{eq:sdf-gap-intro} degrades precisely where a non-convex
        contact forms.
  \item \textbf{Spectral bias.} A coordinate multilayer perceptron learns low frequencies first and high
        frequencies slowly or never; its neural-tangent kernel is Laplace-like, with high-frequency
        power decaying \citep{rahaman2019spectral}. A fixed-capacity neural $\phi_A$ therefore low-passes
        cusps, asperities, and fine detail: it cannot represent a boundary sharper than its capacity
        ceiling, and it smooths the geometry that, on a rough joint, carries the strength.
\end{enumerate}
We measure both weaknesses, rather than merely assert them, in Sections~\ref{sec:verif}
and~\ref{sec:cv7}. The chart introduced next is built to avoid both: it stores one fewer
dimension, and it hands the network the high frequencies it would otherwise fail to learn.

\subsection{Related work}\label{sec:intro-related}

\paragraph{Signed-distance and closest-point contact.}
Reading a gap and a normal from an implicit field is established practice. \citet{kolev2003level} treat
the contact interface as the zero level set of one body's field restricted to the other; the
implicit-level-set material point method of \citet{liu2020implicit} reads the normal gap and contact
normal from a per-body signed-distance representation, and our force machinery reuses its normal-gap
form. The node-to-surface projection of \citet{zavarise1998segment} and \citet{fischer2005frictionless}
obtains a smooth, differentiable correspondence by closest-point projection without explicit facet
intersection, the same smoothness the implicit field provides. It returns only a \emph{single} foot,
however, which is the kinematic limitation the transition map removes (Section~\ref{sec:tmap}).

\paragraph{Neural implicit geometry, spectral bias, and Fourier features.}
Neural signed-distance and occupancy networks \citep{park2019deepsdf} and periodic-activation networks
\citep{sitzmann2020implicit} represent geometry implicitly with a fixed-capacity network, and inherit the
spectral bias characterized by \citet{rahaman2019spectral}: high-frequency detail is learned slowly or
not at all. The remedy we adopt is Fourier-feature input encoding \citep{tancik2020fourier}, which lifts
the input through a fixed bank of sinusoids and flattens the neural-tangent kernel to a tunable cutoff;
the same encoding underlies the positional encoding of neural radiance fields \citep{mildenhall2020nerf}.
We use Fourier features not to sharpen an ambient field but to sharpen a \emph{boundary} chart, which is
one dimension lower and therefore far easier to fit to a prescribed cutoff (Section~\ref{sec:fourier}).

\paragraph{Constraint enforcement.}
The forces that consume the gap and the normal are standard. We use penalty and augmented-Lagrangian
normal contact with regularized Coulomb friction \citep{wriggers2006computational,alart1991mixed,
simo1992augmented}, set against the mortar \citep{puso2004mortar} and Nitsche \citep{chouly2013nitsche}
families and the variational-inequality foundations of \citet{kikuchi1988contact} and
\citet{laursen2002computational}. These schemes are orthogonal to our contribution. Because every one of
them consumes only the gap, the normal, and the normal-force magnitude, the chart and the level set are
interchangeable at the force layer, and the geometry is the only thing we vary.

\paragraph{Rock-joint constitutive laws.}
The validation problem is the direct shear of a rough rock joint, for which a mature constitutive
literature supplies the closed-form anchor and the physical interpretation of the measured numbers. The
bilinear law of \citet{patton1966multiple} gives the apparent friction angle of a regularly serrated
joint as $\phi_b+i$, with $\phi_b$ the basic friction angle and $i$ the asperity inclination; the
empirical model of \citet{barton1977shear} and the review of \citet{barton1990review} extend this to
natural roughness; the asperity-degradation model of \citet{plesha1987constitutive} and the joint element
of \citet{goodman1976methods} supply the interface kinematics for cyclic shear. In our framework the
dilatancy and the apparent friction are not imposed through such a law. They \emph{emerge} from the
resolved asperity geometry, and the Patton law is the reference these emergent quantities are checked
against.

\subsection{Contributions}\label{sec:intro-contributions}

\begin{enumerate}
  \item \textbf{A transition-map contact detector.} We read contact from the boundary-to-boundary chart
        map $\tmap=\map{B}^{-1}\circ\map{A}$ between two bodies of a learned atlas, rather than from an
        ambient field. The detector returns a matched, conservative gap--normal pair, in which the normal
        is the gradient of the gap, so the penalty force is the exact gradient of a potential. A
        boundary scan then enumerates \emph{every} disjoint contact arc, where a closest-point projection
        returns one foot. This is the paper's central method (Section~\ref{sec:tmap}).
  \item \textbf{Fourier-feature coordinate charts with a verify-first protocol.} We instantiate the chart
        in three families: a radial height field $\rho_\theta$ on the sphere of directions, an open
        graph height field $h_\theta$, and a boundary-fitted decoder $D_\theta$ inverted by Newton
        iteration. Each is encoded with a Fourier-feature bank sized to the geometry's resolvable
        bandwidth. We verify every chart against a closed form or a manufactured solution before we use
        it to drive contact (Section~\ref{sec:fourier}).
  \item \textbf{A characterization of when the chart beats the level set.} We show that the radial
        inverse-chart gap is conservative-large but \emph{sign-exact},
        $g_{\mathrm{rad}}=g_\perp/\cos\alpha+O(g^2)$ with $\alpha$ the ray-to-normal angle, so the
        chart's active set matches the true level set's on star-shaped bodies; and we give the
        complementary spectral-bias argument that a $(d{-}1)$-dimensional boundary function fits a
        prescribed cutoff where a $d$-dimensional ambient field cannot
        (Sections~\ref{sec:tmap}--\ref{sec:fourier}).
  \item \textbf{A verification ladder CV-1 through CV-6 and a real-joint validation CV-7.} We place the
        chart and the level set on the same problems: the convex, superposable closed forms
        (CV-1 Hertz, CV-2 Cattaneo--Mindlin, CV-3 Brazilian, CV-4 nine-disc), where they are equivalent
        and the chart buys nothing; the cusped non-convex superformula (CV-5) and the
        self-similar Koch fractal (CV-6), where the chart measurably wins; and the direct shear of a real
        Inada-granite fracture (CV-7), where the resolved asperities drive emergent dilatancy and strength
        (Sections~\ref{sec:verif}--\ref{sec:cv7}).
  \item \textbf{Embedding in chart-FEM and MPM.} The detector is an opt-in geometric oracle: it composes
        with the existing chart-based finite-element and material-point solvers without modifying the
        force laws, the broad-phase culler, or the time integrator (Section~\ref{sec:form}).
\end{enumerate}

% code: benchmarks/contact/cv_numerical/cv7_transition_map_contact.py
A handful of measurements anchor these claims. On the cusped superformula (CV-5) the neural radial chart
attains a gap error of $3.8\times10^{-3}L$ and a $0.42^\circ$ median normal-angle error, against
$8\times10^{-3}L$ and a degraded normal for the neural level set, and it drives the cam dynamics to a
$0.04\%$ trajectory match. On the real Inada joint (CV-7) the height chart reconstructs the surface to
$2.3\,\mu$m where an ambient neural level set reaches only $107\,\mu$m; the level set smooths the mean
asperity angle from $19.4^\circ$ to $12.5^\circ$ and, through the Patton law, under-predicts peak shear
strength by $61\%$. On a band-limited surface carried by a trained decoder the level set under-predicts
emergent dilatancy by $98\%$ and strength by $35\%$, while the transition-map detector in the contact
loop agrees with the analytic active set $98.9\%$ of the time, against $95.8\%$ for the level set
(Sections~\ref{sec:verif}--\ref{sec:cv7}).

\subsection{Notation and organization}\label{sec:intro-notation}

Bodies are labelled $A,B,\dots$; the interior of body $X$ is $\Omega_X\subset\R^d$ ($d=2,3$) and its
boundary $\partial\Omega_X$. A \emph{chart} is a parametrization of a body's boundary,
$\map{X}:\Theta_X\to\partial\Omega_X$; a \emph{level set} is an ambient scalar field $\phi_X:\R^d\to\R$
whose zero set is $\partial\Omega_X$. The transition map between two charted bodies is
$\tmap=\map{B}^{-1}\circ\map{A}$. We write the Macaulay bracket $\pos{z}=\max(0,z)$ and keep the sign
convention $\gn<0\Leftrightarrow$ penetration. The apparent friction coefficient measured from a contact
run is $\muapp$, to be compared against the Patton value $\tan(\phi_b+i)$.

The remainder of the paper is organized as follows. Section~\ref{sec:tmap} defines the
transition map, derives the matched gap--normal pair and the multi-arc enumeration, and states the
$1/\cos\alpha$ radial-gap bound with its active-set equivalence. Section~\ref{sec:fourier} develops the
Fourier-feature charts and the spectral-bias argument that they defeat. Section~\ref{sec:form}
gives the force laws and the embedding in chart-FEM and MPM. Section~\ref{sec:verif} reports the
closed-form ladder CV-1 through CV-6. Section~\ref{sec:cv7} reports the real-joint validation
CV-7. Section~\ref{sec:disc} states the scope, preconditions, and limitations, and
Section~\ref{sec:concl} concludes.
